\section{A finite-cut defect bound}\label{sec:budget}

Throughout this section we assume that there are finitely many prime
geodesic images. Our objective is to compare a sublevel with its
half-translation fixed locus while retaining the information lost by
localization. The comparison has two degrees: the total equivariant
degree upstairs and the degree of a root loop downstairs. An exact-sequence
calculation will bound their discrepancy, summed over all degree cutoffs,
by the number of critical events.

Write
\[
 B_a^l=H_R^l(X_a,M),\qquad
 B_\infty^l=H_R^l(\Lambda,M).
\]
At a regular cut these reflection groups have finite total dimension:
Theorem~\ref{thm:local-control} identifies each reflection crossing with
a finite sum of ordinary one-orientation crossing groups, and there are
only finitely many crossings below a finite cut. All sublevels in the
statements are open. At regular endpoints, the deformations of
Lemmas~\ref{lem:local-strips} and~\ref{lem:local-polygons} identify them, naturally
under inclusion and fixed restriction, with closed regular polygon
models. We use those compact models in the localization argument.

\begin{lemma}[Finite-pair localization]\label{bud:localization}
Let $2a>0$ be regular. Then $a$ is regular, and restriction to the
half-translation fixed locus induces an isomorphism
\[
 H_K^*(X_{2a},M)[v^{-1}]
 \longrightarrow H_K^*(X_{2a}^{\tau},M)[v^{-1}].
\]
It is an actual isomorphism in every sufficiently high degree, with the
threshold allowed to depend on $a$. The same localization assertion
holds for crossing pairs and respects their individual prime labels.

For an ambient crossing at the $m$-th iterate of a prime, its actual
$K$-group has uniformly bounded dimension in every degree, is zero below
$I_m-A$, and maps isomorphically to its stabilized fixed target above
$I_m+A$, after increasing the uniform constant $A$. For odd $m$ that
target is zero; for even $m$ its root labels satisfy
$|l-I_{m/2}|\le A$.
\end{lemma}

\begin{proof}
A critical root doubles to a critical loop, so regularity of $2a$
implies regularity of $a$. Doubling identifies $X_a$ with
$X_{2a}^{\tau}$ as an $R$-space. Since $\tau$ acts trivially there, a
product universal model and the field K\"unneth formula give
\begin{equation}\label{bud:fixed-product}
 H_K^J(X_{2a}^{\tau},M)
 =\bigoplus_{l\le J}v^{J-l}B_a^l.
\end{equation}
The same identification holds for fixed crossing pairs.

Choose a common even polygon model through the cut, and double its
compact regular sublevel along an invariant collar. Denote the resulting
closed smooth $K$-manifold by $\mathcal D$. Inclusion of the original
sublevel and the fold map make its restriction arrow to the fixed locus
a retract of the corresponding arrow for $\mathcal D$. The fixed
submanifold $\mathcal D^\tau$ has an invariant open tubular neighborhood
$O_0$ which retracts equivariantly onto it. The compact complement of
$O_0$ is covered by finitely many invariant open sets $O_1,\ldots,O_s$,
each admitting an equivariant map to $K/H_i$ with $\tau\notin H_i$.
Indeed, use a small stabilizer-invariant neighborhood disjoint from its
translates by other cosets, and then saturate it.

The possible groups $H_i$ are $1$, $R=\langle r\rangle$, and
$\langle\tau r\rangle$.
On each, $U=v(v+w)$ restricts to zero. Thus $U$ lifts to
$H_K^2(\mathcal D,O_i)$. Cup these lifts with a class relative to $O_0$.
The relative cup product for the open cover has target relative to
$\bigcup_{i=0}^sO_i=\mathcal D$, so $U^s$ annihilates
$H_K^*(\mathcal D,O_0)$. The exact sequence of this pair proves
localization at $U$. Passing to the retract proves it on the original
sublevel. The exact sequences for a pair of regular sublevels, and for
the pair relative to $M$, give the corresponding assertions for pairs.

The class $w$ is zero on each ambient local block, because that block
is induced from the $\langle\tau\rangle$-pair of one orientation.
It is therefore nilpotent on each finite-cut relative group by induction
over the crossing exact sequences. The same argument applies to the
fixed reflection groups and to crossing pairs. On these modules
localization at $U$ is equivalent to localization at $v$: after
inverting $v$, the factor $1+w/v$ in $U=v^2(1+w/v)$ is invertible;
after inverting $U$, the element $(v+w)/U$ is an inverse for $v$.

Apply the sign-line Gysin sequence of Proposition~\ref{prop:gysin} to
one orientation of a local block. The total ordinary rank bound and
index support in Theorem~\ref{thm:local-control} imply the claimed
per-degree bound, vanishing below the index window, and that multiplication
by $v$ is an isomorphism above it. The fixed local group has the product
description in \eqref{bud:fixed-product}. Its root support is within a
fixed distance of $I_{m/2}$, and
\[
 |I_m-2I_{m/2}|\le n.
\]
Consequently both source and target have already stabilized above
$I_m+A$, with a uniform enlarged $A$. The localized restriction
isomorphism is then an actual isomorphism there: clear a power of $v$
and invert the intervening, already invertible multiplication maps.
For an odd iterate the fixed block is empty, so its stabilized source
is zero. Labelwise naturality follows from the common full/fixed lens
in Theorem~\ref{thm:local-control}.

Finally, the sign-line Gysin sequence from $K$ to $R$ and the finite
total dimension of $B_{2a}$ make multiplication by $v$ an isomorphism
on the ambient relative group in all sufficiently high degrees.
The fixed group has the same property by \eqref{bud:fixed-product}.
The same clearing argument upgrades localization to actual restriction
in high degrees at this cut.
\end{proof}

For a paired regular cut $(2a,a)$, put
\[
 \mathcal V_a=\bigoplus_l B_a^l,\qquad
 A_J=H_K^J(X_{2a},M).
\]
Dropping the powers $v^{J-l}$ in \eqref{bud:fixed-product} defines the
normalized actual restriction
$\epsilon_J:A_J\to\mathcal V_a$. Set
\[
 E_J=\im\epsilon_J,\qquad t_J=\dim\ker\epsilon_J.
\]
The image $E_J$ lies in labels at most $J$; it need not be graded by
root label, and it need not equal the stabilized target. In these
normalized coordinates, multiplication by $v$ leaves a vector unchanged,
whereas multiplication by $w$ applies the root $w$ to each component and
raises its label by one.

Let $P_j$ denote projection to root labels at most $j$. We record three
defects: the dimension missing from the projected low root degrees,
the dimension of actual-image vectors supported entirely above the cutoff,
and the dimension of actual classes invisible on the fixed locus.
They are, respectively, $D_{j,q}$, $H_{j,J}$, and $t_J$.
The potential also includes the preceding-degree deficit $D_{j,2j-1}$;
this compensates for the degree shift in the connecting homomorphism.
Define
\begin{equation}\label{bud:potential}
 \begin{split}
 D_{j,q}&=\dim P_j\mathcal V_a-\dim P_jE_q,\\
 H_{j,J}&=\dim(E_J\cap\ker P_j),\qquad J=2j,\\
 \Phi_j(a)&=D_{j,2j}+D_{j,2j-1}+H_{j,2j}+t_{2j},
 \qquad j\ge1.
 \end{split}
\end{equation}
All four terms are nonnegative. Lemma~\ref{bud:localization} shows that
they vanish for sufficiently large $j$ at any fixed cut.

\begin{theorem}[Finite-cut defect bound]\label{thm:finite-cut-budget}
There is a constant $C$ such that, whenever $2a>0$ is regular,
\begin{equation}\label{bud:linear-budget}
 \sum_{j\ge1}\Phi_j(a)\le C(1+a).
\end{equation}
\end{theorem}

\begin{proof}
Grow the ambient cut through its critical events, using primes for the
groups just after an event. Let $C^q$ be its actual local $K$-group and
$\mathcal W$ its total fixed root reflection group. The comparison
maps fit into the following diagram, whose rows are exact:
\begin{equation}\label{bud:crossing-diagram}
 \begin{tikzcd}[column sep=small]
 A_{q-1}\arrow[r]\arrow[d,"\epsilon_{q-1}"'] &
 C^q\arrow[r]\arrow[d,"\lambda_q"'] &
 A'_q\arrow[r]\arrow[d,"\epsilon'_q"'] &
 A_q\arrow[r]\arrow[d,"\epsilon_q"'] &
 C^{q+1}\arrow[d,"\lambda_{q+1}"]\\
 \mathcal V\arrow[r,"\delta"'] &
 \mathcal W\arrow[r,"f"'] &
 \mathcal V'\arrow[r,"i"'] &
 \mathcal V\arrow[r,"\delta"'] &\mathcal W.
 \end{tikzcd}
\end{equation}
The bottom row is the fixed exact sequence after inverting $v$ and
normalizing powers consistently in neighboring degrees. Its maps $f,i$
preserve root labels, and $\delta$ raises a root label by one. In
particular, the repeated terms denote the total graded groups in an
exact sequence, not an assertion of ungraded periodicity.

Fix $j\ge1$ and put $J=2j$. Call this cutoff \emph{far} from the event
if every participating iterate satisfies
$|I_m-2j|>C_0$, for a sufficiently large fixed $C_0$.
Lemma~\ref{bud:localization} and the index comparison then give
\begin{equation}\label{bud:far-local}
 \mathcal W^j=\mathcal W^{j+1}=0,
 \qquad
 \lambda_{J-1},\lambda_J,\lambda_{J+1}
 \text{ injective with image }P_j\mathcal W.
\end{equation}
Indeed, the blocks below the window are already in their actual stable
range and have all root labels below $j$; those above it have zero
actual source in these three degrees and no low root labels.

Write $P=P_j$, $L=P\mathcal V$, $L'=P\mathcal V'$, and
\[
 A_0=f(P\mathcal W),\qquad
 r=\rank(\delta|_L),\qquad
 r_q=\rank(\delta|_{PE_q})\quad(q=J-1,J).
\]
There is no connecting term across the projection edge, so low-label
exactness gives
\begin{equation}\label{bud:low-dimensions}
 \dim P\mathcal W=\dim A_0+r,
 \qquad \dim L'-\dim L=\dim A_0-r.
\end{equation}

We next justify the projected-image calculation, including mixed low
and high components. For $z\in E_q$, commutativity and
\eqref{bud:far-local} place $\delta z$ in $P\mathcal W$.
Also $\delta(Pz)$ is low, since $\mathcal W^{j+1}=0$, while
$\delta((1-P)z)$ has labels greater than $j$. These disjoint label
ranges imply
\[
 \delta((1-P)z)=0,\qquad \delta z=\delta(Pz).
\]
Given $x\in PE_q\cap\ker\delta$, choose $z\in E_q$ projecting to $x$
and an actual representative in $A_q$. Its actual connecting class
restricts to zero. The target $\lambda_{q+1}$ is injective, so actual
exactness lifts the representative to $A'_q$. Projecting its fixed
restriction proves
\[
 i(PE'_q)=PE_q\cap\ker\delta.
\]
The kernel of this map is exactly $A_0$: stable exactness gives the
inclusion in $A_0$, and every element of $P\mathcal W$ has an actual
local lift. Therefore
\begin{equation}\label{bud:deficit-change}
 D'_{j,q}-D_{j,q}=-(r-r_q),\qquad q=J-1,J.
\end{equation}
The actual connecting ranks are precisely $r_q$, again by target
injectivity. Hence
\[
 \dim A'_J-\dim A_J=\dim P\mathcal W-r_{J-1}-r_J.
\]
Using $H_{j,J}+t_J=\dim A_J-\dim PE_J$ and
\eqref{bud:low-dimensions} gives
\[
 (H_{j,J}+t_J)'-(H_{j,J}+t_J)=r-r_{J-1}.
\]
Together with \eqref{bud:deficit-change}, this proves the exact change
\begin{equation}\label{bud:far-change}
 \Phi'_j-\Phi_j=-(r-r_J)\le0.
\end{equation}
The preceding-degree deficit cancels the possible increase in the
high-image and kernel terms. This is why all four terms were retained.

For a cutoff which is not far, local restriction need not be injective.
Nevertheless the change is uniformly bounded. Choose $M_0$ bounding
$\dim\mathcal W$ and all the relevant actual local dimensions. Since
$i$ preserves labels, $i(PE'_q)\subset PE_q$. Its kernel has dimension
at most $M_0$ by stable exactness. Its cokernel is bounded using the
surjections
\[
 A_q/\im A'_q\ \longrightarrow\ E_q/iE'_q
 \ \longrightarrow\ PE_q/i(PE'_q).
\]
The first source has dimension at most $\dim C^{q+1}\le M_0$.
Thus $|\dim PE'_q-\dim PE_q|\le M_0$.
The same kernel and cokernel argument bounds
$|\dim L'-\dim L|$ by $M_0$, and actual exactness bounds
$|\dim A'_J-\dim A_J|$ by $M_0$. Since
\[
 \Phi_j=2\dim L-2\dim PE_J-\dim PE_{J-1}+\dim A_J,
\]
we obtain $|\Phi'_j-\Phi_j|\le6M_0$.

Each prime contributes at most one iterate to a given positive event.
The finite prime list therefore bounds the number of participating
blocks, and only a bounded number of integers $j$ lie in their fixed
index windows. Thus every event increases the sum of the potentials by
at most a fixed constant. The number of events below $2a$ is at most
$\sum_i\lfloor2a/\ell_i\rfloor$. The potentials start at zero below
the first positive event and are unchanged on no-critical strips.
Their finite support at each cut permits summation of
\eqref{bud:far-change} and the exceptional bounds, proving
\eqref{bud:linear-budget}.
\end{proof}

The estimate controls all degree cutoffs at once. In the next section,
a single persistent extra reflection class will produce quadratically
many units of this linear budget. Degree zero, which is not included in
\eqref{bud:potential}, will be bounded separately by ordinary death.
